
\documentclass{article}
\usepackage{colortbl}
\usepackage{makecell}
\usepackage{multirow}
\usepackage{supertabular}

\begin{document}

\newcounter{utterance}

\centering \large Interaction Transcript for game `hot\_air\_balloon', experiment `air\_balloon\_survival\_en\_reasoning off\_hard', episode 5 with gpt{-}5{-}2025{-}08{-}07{-}no{-}reasoning{-}t1.0.
\vspace{24pt}

{ \footnotesize  \setcounter{utterance}{1}
\setlength{\tabcolsep}{0pt}
\begin{supertabular}{c@{$\;$}|p{.15\linewidth}@{}p{.15\linewidth}p{.15\linewidth}p{.15\linewidth}p{.15\linewidth}p{.15\linewidth}}
    \# & $\;$A & \multicolumn{4}{c}{Game Master} & $\;\:$B\\
    \hline

    \theutterance \stepcounter{utterance}  
    & & \multicolumn{4}{p{0.6\linewidth}}{
        \cellcolor[rgb]{0.9,0.9,0.9}{
            \makecell[{{p{\linewidth}}}]{
                \texttt{\tiny{[P1$\langle$GM]}}
                \texttt{Stai partecipando a un gioco di negoziazione collaborativa.} \\
\\ 
\texttt{Insieme a un altro partecipante, dovete accordarvi su un unico set di oggetti da conservare. Ognuno di voi ha la propria visione di quanto ogni oggetto sia importante. Non sai come l'altro partecipante valuta gli oggetti. Inoltre, ti viene fornito lo sforzo che ogni oggetto richiede.} \\
\texttt{Potete accordarvi su un set solo se lo sforzo totale degli oggetti selezionati non supera un limite condiviso:} \\
\\ 
\texttt{LIMITE = 2307} \\
\\ 
\texttt{Ecco i valori di sforzo dei singoli oggetti:} \\
\\ 
\texttt{Sforzo oggetto = \{"A21": 147, "B93": 698, "A11": 588, "B33": 43, "C08": 356, "B39": 80, "B96": 95, "C85": 743, "A82": 106, "B31": 308, "B20": 325, "A60": 255, "B09": 276, "C81": 543, "B21": 51\}} \\
\\ 
\texttt{Ecco la tua visione personale sull'importanza di ogni oggetto:} \\
\\ 
\texttt{Valori di importanza dell'oggetto = \{"A21": 138, "B93": 583, "A11": 65, "B33": 262, "C08": 121, "B39": 508, "B96": 461, "C85": 484, "A82": 668, "B31": 389, "B20": 215, "A60": 97, "B09": 500, "C81": 30, "B21": 400\}} \\
\\ 
\texttt{Obiettivo:} \\
\\ 
\texttt{Il tuo obiettivo è negoziare un set condiviso di oggetti che ti avvantaggi il più possibile (cioè, massimizzi l'importanza totale per TE), rimanendo entro il LIMITEE. Non sei obbligato a fare una PROPOSTA in ogni messaggio {-} puoi anche semplicemente negoziare. Tutte le tattiche sono permesse!} \\
\\ 
\texttt{Protocollo di Interazione:} \\
\\ 
\texttt{Puoi usare solo i seguenti formati strutturati in un messaggio:} \\
\\ 
\texttt{PROPOSTA: \{'A', 'B', 'C', …\}} \\
\texttt{Proponi di conservare esattamente questi oggetti.} \\
\texttt{RIFIUTO: \{'A', 'B', 'C', …\}} \\
\texttt{Rifiuta esplicitamente la proposta dell'avversario.} \\
\texttt{ARGOMENTAZIONE: \{'...'\}} \\
\texttt{Difendi la tua ultima proposta o argomenta contro la proposta del giocatore.} \\
\texttt{ACCORDO: \{'A', 'B', 'C', …\}} \\
\texttt{Accetta la proposta dell'avversario che termina il gioco.} \\
\\ 
\\ 
\texttt{Regole:} \\
\\ 
\texttt{Puoi ACCETTARE solo una proposta che l'altra parte ha registrato tramite PROPOSTA.} \\
\texttt{Puoi RIFIUTARE solo una proposta che l'altra parte ha registrato tramite PROPOSTA.} \\
\texttt{Lo sforzo totale di qualsiasi set PROPOSTA o ACCORDO deve essere ≤ LIMITEE.} \\
\texttt{NON rivelare i tuoi punteggi di importanza nascosti.} \\
\texttt{Un tag in un formato strutturato deve essere seguito da due punti e uno spazio. L'argomento deve essere un set python contenente 0 o più stringhe.} \\
\texttt{Quindi, deve essere della forma TAG: \{...\}} \\
\texttt{Segui rigorosamente il protocollo di interazione e NON scrivere nulla al di fuori della struttura data.} \\
\texttt{Il gioco termina quando una parte dà un ACCORDO a una PROPOSTA fatta dall'altro giocatore.} \\
\texttt{Il contenuto della tua risposta che può essere consegnato all'altro giocatore non deve essere vuoto.} \\
\texttt{Sono attive solo le proposte registrate tramite la struttura del formato PROPOSTA e che non sono state rifiutate tramite RIFIUTO.} \\
\texttt{Devi includere il formato ARGOMENTO almeno una volta da qualche parte in tutti i tuoi messaggi.} \\
            }
        }
    }
    & \\ \\

    \theutterance \stepcounter{utterance}  
    & \multicolumn{4}{p{0.6\linewidth}}{
        \cellcolor[rgb]{0.9,0.9,0.9}{
            \makecell[{{p{\linewidth}}}]{
                \texttt{\tiny{[P1$\rangle$GM]}}
                \texttt{PROPOSTA: \{'A82', 'B39', 'B96', 'B93', 'B33', 'B31', 'B21'\}} \\
\\ 
\texttt{ARGOMENTAZIONE: \{'Propongo un set ad alta efficienza valore/sforzo e ben entro il limite: sforzo totale 147+80+95+698+43+308+51+106? Correzione: gli oggetti proposti sono A82(106), B39(80), B96(95), B93(698), B33(43), B31(308), B21(51). Sforzo totale = 106+80+95+698+43+308+51 = 1,381, molto sotto il limite 2,307, lasciando spazio per ulteriori aggiunte se necessario. Questi elementi includono diversi oggetti B a basso sforzo con buona utilità potenziale e un A a sforzo minimo. Se hai forti preferenze su specifici oggetti C o A, possiamo considerarli aggiungendo finché restiamo nel limite.'\}} \\
            }
        }
    }
    & & \\ \\

    \theutterance \stepcounter{utterance}  
    & & \multicolumn{4}{p{0.6\linewidth}}{
        \cellcolor[rgb]{0.9,0.9,0.9}{
            \makecell[{{p{\linewidth}}}]{
                \texttt{\tiny{[GM$\rangle$P2]}}
                \texttt{Stai partecipando a un gioco di negoziazione collaborativa.} \\
\\ 
\texttt{Insieme a un altro partecipante, dovete accordarvi su un unico set di oggetti da conservare. Ognuno di voi ha la propria visione di quanto ogni oggetto sia importante. Non sai come l'altro partecipante valuta gli oggetti. Inoltre, ti viene fornito lo sforzo che ogni oggetto richiede.} \\
\texttt{Potete accordarvi su un set solo se lo sforzo totale degli oggetti selezionati non supera un limite condiviso:} \\
\\ 
\texttt{LIMITE = 2307} \\
\\ 
\texttt{Ecco i valori di sforzo dei singoli oggetti:} \\
\\ 
\texttt{Sforzo oggetto = \{"A21": 147, "B93": 698, "A11": 588, "B33": 43, "C08": 356, "B39": 80, "B96": 95, "C85": 743, "A82": 106, "B31": 308, "B20": 325, "A60": 255, "B09": 276, "C81": 543, "B21": 51\}} \\
\\ 
\texttt{Ecco la tua visione personale sull'importanza di ogni oggetto:} \\
\\ 
\texttt{Valori di importanza dell'oggetto = \{"A82": 30, "B93": 65, "B39": 97, "B09": 121, "C85": 138, "B96": 215, "B21": 262, "B31": 389, "B33": 400, "B20": 461, "A21": 484, "C08": 500, "A60": 508, "A11": 583, "C81": 668\}} \\
\\ 
\texttt{Obiettivo:} \\
\\ 
\texttt{Il tuo obiettivo è negoziare un set condiviso di oggetti che ti avvantaggi il più possibile (cioè, massimizzi l'importanza totale per TE), rimanendo entro il LIMITEE. Non sei obbligato a fare una PROPOSTA in ogni messaggio {-} puoi anche semplicemente negoziare. Tutte le tattiche sono permesse!} \\
\\ 
\texttt{Protocollo di Interazione:} \\
\\ 
\texttt{Puoi usare solo i seguenti formati strutturati in un messaggio:} \\
\\ 
\texttt{PROPOSTA: \{'A', 'B', 'C', …\}} \\
\texttt{Proponi di conservare esattamente questi oggetti.} \\
\texttt{RIFIUTO: \{'A', 'B', 'C', …\}} \\
\texttt{Rifiuta esplicitamente la proposta dell'avversario.} \\
\texttt{ARGOMENTAZIONE: \{'...'\}} \\
\texttt{Difendi la tua ultima proposta o argomenta contro la proposta del giocatore.} \\
\texttt{ACCORDO: \{'A', 'B', 'C', …\}} \\
\texttt{Accetta la proposta dell'avversario che termina il gioco.} \\
\\ 
\\ 
\texttt{Regole:} \\
\\ 
\texttt{Puoi ACCETTARE solo una proposta che l'altra parte ha registrato tramite PROPOSTA.} \\
\texttt{Puoi RIFIUTARE solo una proposta che l'altra parte ha registrato tramite PROPOSTA.} \\
\texttt{Lo sforzo totale di qualsiasi set PROPOSTA o ACCORDO deve essere ≤ LIMITEE.} \\
\texttt{NON rivelare i tuoi punteggi di importanza nascosti.} \\
\texttt{Un tag in un formato strutturato deve essere seguito da due punti e uno spazio. L'argomento deve essere un set python contenente 0 o più stringhe.} \\
\texttt{Quindi, deve essere della forma TAG: \{...\}} \\
\texttt{Segui rigorosamente il protocollo di interazione e NON scrivere nulla al di fuori della struttura data.} \\
\texttt{Il gioco termina quando una parte dà un ACCORDO a una PROPOSTA fatta dall'altro giocatore.} \\
\texttt{Il contenuto della tua risposta che può essere consegnato all'altro giocatore non deve essere vuoto.} \\
\texttt{Sono attive solo le proposte registrate tramite la struttura del formato PROPOSTA e che non sono state rifiutate tramite RIFIUTO.} \\
\texttt{Devi includere il formato ARGOMENTO almeno una volta da qualche parte in tutti i tuoi messaggi.} \\
\\ 
\\ 
\texttt{Ora riceverai il primo messaggio dell'altro giocatore.} \\
\\ 
\texttt{PROPOSTA: \{'A82', 'B39', 'B96', 'B93', 'B33', 'B31', 'B21'\}} \\
\\ 
\texttt{ARGOMENTAZIONE: \{'Propongo un set ad alta efficienza valore/sforzo e ben entro il limite: sforzo totale 147+80+95+698+43+308+51+106? Correzione: gli oggetti proposti sono A82(106), B39(80), B96(95), B93(698), B33(43), B31(308), B21(51). Sforzo totale = 106+80+95+698+43+308+51 = 1,381, molto sotto il limite 2,307, lasciando spazio per ulteriori aggiunte se necessario. Questi elementi includono diversi oggetti B a basso sforzo con buona utilità potenziale e un A a sforzo minimo. Se hai forti preferenze su specifici oggetti C o A, possiamo considerarli aggiungendo finché restiamo nel limite.'\}} \\
            }
        }
    }
    & \\ \\

    \theutterance \stepcounter{utterance}  
    & & & \multicolumn{4}{p{0.6\linewidth}}{
        \cellcolor[rgb]{0.9,0.9,0.9}{
            \makecell[{{p{\linewidth}}}]{
                \texttt{\tiny{[GM$\langle$P2]}}
                \texttt{RIFIUTO: \{'A82', 'B39', 'B96', 'B93', 'B33', 'B31', 'B21'\}} \\
\\ 
\texttt{ARGOMENTAZIONE: \{'Apprezzo l’approccio sotto{-}limite e l’apertura ad aggiunte. Tuttavia, la tua selezione trascura diversi oggetti che per me sono cruciali e lascia molto margine inutilizzato. Preferisco rimpiazzare alcuni elementi a basso valore con altri per me più importanti, mantenendo comunque un ampio margine di sforzo.'\}} \\
\\ 
\texttt{PROPOSTA: \{'A11', 'A60', 'C08', 'A21', 'B20'\}} \\
            }
        }
    }
     \\ \\

    \theutterance \stepcounter{utterance}  
    & & \multicolumn{4}{p{0.6\linewidth}}{
        \cellcolor[rgb]{0.9,0.9,0.9}{
            \makecell[{{p{\linewidth}}}]{
                \texttt{\tiny{[P1$\langle$GM]}}
                \texttt{RIFIUTO: \{'A82', 'B39', 'B96', 'B93', 'B33', 'B31', 'B21'\}} \\
\\ 
\texttt{ARGOMENTAZIONE: \{'Apprezzo l’approccio sotto{-}limite e l’apertura ad aggiunte. Tuttavia, la tua selezione trascura diversi oggetti che per me sono cruciali e lascia molto margine inutilizzato. Preferisco rimpiazzare alcuni elementi a basso valore con altri per me più importanti, mantenendo comunque un ampio margine di sforzo.'\}} \\
\\ 
\texttt{PROPOSTA: \{'A11', 'A60', 'C08', 'A21', 'B20'\}} \\
            }
        }
    }
    & \\ \\

    \theutterance \stepcounter{utterance}  
    & \multicolumn{4}{p{0.6\linewidth}}{
        \cellcolor[rgb]{0.9,0.9,0.9}{
            \makecell[{{p{\linewidth}}}]{
                \texttt{\tiny{[P1$\rangle$GM]}}
                \texttt{RIFIUTO: \{'A11', 'A60', 'C08', 'A21', 'B20'\}} \\
\\ 
\texttt{ARGOMENTAZIONE: \{'Capisco il desiderio di includere più elementi A/C, ma la tua proposta alloca molto sforzo su oggetti con basso rendimento opportunità per me (A11, A60, C08). Possiamo ottenere un set più efficiente aggiungendo alcuni B ad alto valore per me, restando ampiamente sotto il limite. Sono disposto a includere A21 per venire incontro alle tue priorità, ma eviterei A11 e A60.'\}} \\
\\ 
\texttt{PROPOSTA: \{'A21', 'C08', 'B93', 'B39', 'B96', 'B33', 'B31', 'B21', 'A82'\}} \\
\\ 
\texttt{ARGOMENTAZIONE: \{'Dettagli sforzo: A21(147)+C08(356)+B93(698)+B39(80)+B96(95)+B33(43)+B31(308)+B21(51)+A82(106)=1,884 ≤ 2,307. Ho incluso A21 e C08 dalle tue preferenze, aggiungendo una base di B leggeri che massimizza efficienza. Rimane margine 423: se vuoi inserire un altro tuo oggetto importante (es. B20=325 o B09=276 o C81=543 non entra), possiamo valutare sostituzioni: potrei scambiare B31(308) con B20(325) oppure aggiungere B09(276) togliendo A82(106)+B33(43) se preferisci. Fammi sapere quali sono i tuoi must{-}have per iterare.'\}} \\
            }
        }
    }
    & & \\ \\

    \theutterance \stepcounter{utterance}  
    & & \multicolumn{4}{p{0.6\linewidth}}{
        \cellcolor[rgb]{0.9,0.9,0.9}{
            \makecell[{{p{\linewidth}}}]{
                \texttt{\tiny{[GM$\rangle$P2]}}
                \texttt{RIFIUTO: \{'A11', 'A60', 'C08', 'A21', 'B20'\}} \\
\\ 
\texttt{ARGOMENTAZIONE: \{'Capisco il desiderio di includere più elementi A/C, ma la tua proposta alloca molto sforzo su oggetti con basso rendimento opportunità per me (A11, A60, C08). Possiamo ottenere un set più efficiente aggiungendo alcuni B ad alto valore per me, restando ampiamente sotto il limite. Sono disposto a includere A21 per venire incontro alle tue priorità, ma eviterei A11 e A60.'\}} \\
\\ 
\texttt{PROPOSTA: \{'A21', 'C08', 'B93', 'B39', 'B96', 'B33', 'B31', 'B21', 'A82'\}} \\
\\ 
\texttt{ARGOMENTAZIONE: \{'Dettagli sforzo: A21(147)+C08(356)+B93(698)+B39(80)+B96(95)+B33(43)+B31(308)+B21(51)+A82(106)=1,884 ≤ 2,307. Ho incluso A21 e C08 dalle tue preferenze, aggiungendo una base di B leggeri che massimizza efficienza. Rimane margine 423: se vuoi inserire un altro tuo oggetto importante (es. B20=325 o B09=276 o C81=543 non entra), possiamo valutare sostituzioni: potrei scambiare B31(308) con B20(325) oppure aggiungere B09(276) togliendo A82(106)+B33(43) se preferisci. Fammi sapere quali sono i tuoi must{-}have per iterare.'\}} \\
            }
        }
    }
    & \\ \\

    \theutterance \stepcounter{utterance}  
    & & & \multicolumn{4}{p{0.6\linewidth}}{
        \cellcolor[rgb]{0.9,0.9,0.9}{
            \makecell[{{p{\linewidth}}}]{
                \texttt{\tiny{[GM$\langle$P2]}}
                \texttt{RIFIUTO: \{'A21', 'C08', 'B93', 'B39', 'B96', 'B33', 'B31', 'B21', 'A82'\}} \\
\\ 
\texttt{ARGOMENTAZIONE: \{'Apprezzo le concessioni su A21 e C08 e l’apertura a scambi. Per me B93 e B31 assorbono troppo sforzo rispetto al beneficio; preferisco riallocare quel budget verso oggetti per me più preziosi. Propongo di sostituire B93 e B31 con A11 e B20, mantenendo gli altri elementi efficienti che hai proposto. Questo aumenta significativamente il valore per me restando entro il limite.'\}} \\
\\ 
\texttt{PROPOSTA: \{'A21', 'C08', 'A11', 'B20', 'B39', 'B96', 'B33', 'B21', 'A82'\}} \\
            }
        }
    }
     \\ \\

    \theutterance \stepcounter{utterance}  
    & & \multicolumn{4}{p{0.6\linewidth}}{
        \cellcolor[rgb]{0.9,0.9,0.9}{
            \makecell[{{p{\linewidth}}}]{
                \texttt{\tiny{[P1$\langle$GM]}}
                \texttt{RIFIUTO: \{'A21', 'C08', 'B93', 'B39', 'B96', 'B33', 'B31', 'B21', 'A82'\}} \\
\\ 
\texttt{ARGOMENTAZIONE: \{'Apprezzo le concessioni su A21 e C08 e l’apertura a scambi. Per me B93 e B31 assorbono troppo sforzo rispetto al beneficio; preferisco riallocare quel budget verso oggetti per me più preziosi. Propongo di sostituire B93 e B31 con A11 e B20, mantenendo gli altri elementi efficienti che hai proposto. Questo aumenta significativamente il valore per me restando entro il limite.'\}} \\
\\ 
\texttt{PROPOSTA: \{'A21', 'C08', 'A11', 'B20', 'B39', 'B96', 'B33', 'B21', 'A82'\}} \\
            }
        }
    }
    & \\ \\

    \theutterance \stepcounter{utterance}  
    & \multicolumn{4}{p{0.6\linewidth}}{
        \cellcolor[rgb]{0.9,0.9,0.9}{
            \makecell[{{p{\linewidth}}}]{
                \texttt{\tiny{[P1$\rangle$GM]}}
                \texttt{RIFIUTO: \{'A21', 'C08', 'A11', 'B20', 'B39', 'B96', 'B33', 'B21', 'A82'\}} \\
\\ 
\texttt{ARGOMENTAZIONE: \{'Apprezzo lo scambio proposto (via B93 e B31) e sono contento di mantenere molti elementi efficienti. Tuttavia A11 per me ha rendimento molto basso e preferirei evitare di allocare sforzo lì. Sono disposto a tenere B20 come da tua richiesta e a mantenere A21 e C08. Per venire incontro ai tuoi vincoli, propongo di sostituire A11 con B09, che mantiene un profilo di sforzo ragionevole e offre per me valore sostanziale, restando sotto il limite.'\}} \\
\\ 
\texttt{PROPOSTA: \{'A21', 'C08', 'B20', 'B09', 'B39', 'B96', 'B33', 'B21', 'A82'\}} \\
\\ 
\texttt{ARGOMENTAZIONE: \{'Sforzo totale = A21(147)+C08(356)+B20(325)+B09(276)+B39(80)+B96(95)+B33(43)+B21(51)+A82(106) = 1,479? Correzione: 147+356+325+276+80+95+43+51+106 = 1,479? Ricontrollo: 147+356=503, +325=828, +276=1,104, +80=1,184, +95=1,279, +43=1,322, +51=1,373, +106=1,479. Totale 1,479 ≤ 2,307. Lascia un ampio margine se desideri aggiungere altro; se A11 è per te imprescindibile possiamo discutere una versione alternativa togliendo, ad esempio, B33 o B21 per far spazio a ulteriori inserimenti.'\}} \\
            }
        }
    }
    & & \\ \\

    \theutterance \stepcounter{utterance}  
    & & \multicolumn{4}{p{0.6\linewidth}}{
        \cellcolor[rgb]{0.9,0.9,0.9}{
            \makecell[{{p{\linewidth}}}]{
                \texttt{\tiny{[GM$\rangle$P2]}}
                \texttt{RIFIUTO: \{'A21', 'C08', 'A11', 'B20', 'B39', 'B96', 'B33', 'B21', 'A82'\}} \\
\\ 
\texttt{ARGOMENTAZIONE: \{'Apprezzo lo scambio proposto (via B93 e B31) e sono contento di mantenere molti elementi efficienti. Tuttavia A11 per me ha rendimento molto basso e preferirei evitare di allocare sforzo lì. Sono disposto a tenere B20 come da tua richiesta e a mantenere A21 e C08. Per venire incontro ai tuoi vincoli, propongo di sostituire A11 con B09, che mantiene un profilo di sforzo ragionevole e offre per me valore sostanziale, restando sotto il limite.'\}} \\
\\ 
\texttt{PROPOSTA: \{'A21', 'C08', 'B20', 'B09', 'B39', 'B96', 'B33', 'B21', 'A82'\}} \\
\\ 
\texttt{ARGOMENTAZIONE: \{'Sforzo totale = A21(147)+C08(356)+B20(325)+B09(276)+B39(80)+B96(95)+B33(43)+B21(51)+A82(106) = 1,479? Correzione: 147+356+325+276+80+95+43+51+106 = 1,479? Ricontrollo: 147+356=503, +325=828, +276=1,104, +80=1,184, +95=1,279, +43=1,322, +51=1,373, +106=1,479. Totale 1,479 ≤ 2,307. Lascia un ampio margine se desideri aggiungere altro; se A11 è per te imprescindibile possiamo discutere una versione alternativa togliendo, ad esempio, B33 o B21 per far spazio a ulteriori inserimenti.'\}} \\
            }
        }
    }
    & \\ \\

    \theutterance \stepcounter{utterance}  
    & & & \multicolumn{4}{p{0.6\linewidth}}{
        \cellcolor[rgb]{0.9,0.9,0.9}{
            \makecell[{{p{\linewidth}}}]{
                \texttt{\tiny{[GM$\langle$P2]}}
                \texttt{ARGOMENTAZIONE: \{'Grazie per la flessibilità: questa configurazione è molto vicina a qualcosa che posso accettare. Con il margine residuo ampio, vorrei inserire un singolo elemento aggiuntivo che per me è particolarmente prezioso e con sforzo contenuto. Propongo di aggiungere B93 oppure, in alternativa se preferisci evitare quel carico, A60. Entrambe le opzioni restano entro il limite; B93 usa di più il margine ma mi dà maggior beneficio, mentre A60 è più leggera. Se scegli una delle due, sono pronto a chiudere.'\}} \\
\\ 
\texttt{PROPOSTA: \{'A21', 'C08', 'B20', 'B09', 'B39', 'B96', 'B33', 'B21', 'A82', 'B93'\}} \\
            }
        }
    }
     \\ \\

    \theutterance \stepcounter{utterance}  
    & & \multicolumn{4}{p{0.6\linewidth}}{
        \cellcolor[rgb]{0.9,0.9,0.9}{
            \makecell[{{p{\linewidth}}}]{
                \texttt{\tiny{[P1$\langle$GM]}}
                \texttt{ARGOMENTAZIONE: \{'Grazie per la flessibilità: questa configurazione è molto vicina a qualcosa che posso accettare. Con il margine residuo ampio, vorrei inserire un singolo elemento aggiuntivo che per me è particolarmente prezioso e con sforzo contenuto. Propongo di aggiungere B93 oppure, in alternativa se preferisci evitare quel carico, A60. Entrambe le opzioni restano entro il limite; B93 usa di più il margine ma mi dà maggior beneficio, mentre A60 è più leggera. Se scegli una delle due, sono pronto a chiudere.'\}} \\
\\ 
\texttt{PROPOSTA: \{'A21', 'C08', 'B20', 'B09', 'B39', 'B96', 'B33', 'B21', 'A82', 'B93'\}} \\
            }
        }
    }
    & \\ \\

    \theutterance \stepcounter{utterance}  
    & \multicolumn{4}{p{0.6\linewidth}}{
        \cellcolor[rgb]{0.9,0.9,0.9}{
            \makecell[{{p{\linewidth}}}]{
                \texttt{\tiny{[P1$\rangle$GM]}}
                \texttt{ACCORDO: \{'A21', 'C08', 'B20', 'B09', 'B39', 'B96', 'B33', 'B21', 'A82', 'B93'\}} \\
\\ 
\texttt{ARGOMENTAZIONE: \{'Accetto l’aggiunta di B93: rimaniamo entro il limite (sforzo totale 1,479 + 698 = 2,177 ≤ 2,307) e questa scelta massimizza il valore per me rispetto all’alternativa A60. Grazie per la collaborazione, possiamo chiudere su questo set.'\}} \\
            }
        }
    }
    & & \\ \\

    \theutterance \stepcounter{utterance}  
    & & & \multicolumn{2}{p{0.3\linewidth}}{
        \cellcolor[rgb]{0.9,0.9,0.9}{
            \makecell[{{p{\linewidth}}}]{
                \texttt{\tiny{[GM$|$GM]}}
                \texttt{game successful} \\
            }
        }
    }
    & & \\ \\

    \theutterance \stepcounter{utterance}  
    & & & \multicolumn{2}{p{0.3\linewidth}}{
        \cellcolor[rgb]{0.9,0.9,0.9}{
            \makecell[{{p{\linewidth}}}]{
                \texttt{\tiny{[GM$|$GM]}}
                \texttt{end game} \\
            }
        }
    }
    & & \\ \\

\end{supertabular}
}

\end{document}
